\section{A genuine citation can still support the wrong answer}
\label{sec:legal-ai-source-support}

The legal-AI evaluation literature adds a failure that a source-preservation
system alone cannot detect. An answer may link to an authentic document and
still misstate its consequence, use an inapplicable jurisdiction, or rely on a
position that has changed. The review therefore needs two questions: whether
the stated proposition is correct, and whether the cited authority supports it
for the activity and time under assessment.

\citet[sections 4--5 and appendix C]{legalrag2024} distinguish factual
correctness from legal groundedness when evaluating legal research tools. Their
2024 study used 202 queries, recorded product responses and manually applied
a published coding rubric. A source can be real yet fail to support the
proposition attributed to it. Their treatment of partially correct answers and
refusals also shows why a label must be read with its scoring definition. The
study's products, dates and United States legal questions do not estimate
LegalMath's error rate on Hong Kong circulars. Its useful contribution here is
the separation of answer correctness, citation support and incomplete answers.

\begin{figure}[H]
\centering
\TeachingCompare{Authentic document}{Supported proposition}{Applicable authority and date}
\caption{Three checks are needed before a citation can support a regulatory control.}
\label{fig:citation-three-checks}
\end{figure}

Consider a synthetic drafting error. A candidate states that an SFC request is
always required before a particular action, citing a paragraph that requires a
response \emph{when} the SFC requests it. The document exists and the quoted
words are accurate, but the candidate has changed the trigger. The audit should
retain the candidate's precise claim, the relevant sentence and its surrounding
conditions, and the reason the implication is unsupported. A quotation is useful
because another reviewer can inspect that comparison; quotation marks themselves
do not establish that the inference is sound.

\citet[sections 4.2--4.3 and appendix C]{legalfictions2024} distinguish
reference-based checks from a method that detects contradictions between model
responses. They inspect United States case-law tasks and validate the
contradiction labelling against human recoding of a sample. The distinction
matters to an ensemble: detecting a contradiction identifies an inconsistent
pair, while agreement leaves shared error possible. The workbench consequently
needs source support even when repeated answers are consistent. The paper's
published error estimates and its model-based labelling are not a substitute
for an independent Hong Kong reference process.

A citation audit must also record the edition actually read. A preprint, a
conference version, a later journal article and a changing software repository
may support different statements. If a later edition changes a theorem or a
benchmark denominator, retaining only the newest file destroys the explanation
of an earlier claim. The appropriate record binds the document bytes, section or
page, a short exact quotation, the manuscript proposition and the support
judgment. The full retained document allows examination beyond the short quote.

\section{Where the literature still leaves this product exposed}
\label{sec:remaining-literature-gaps}

The collection is strongest on executable legal languages, formal comparisons
and candidate-generation methods. Those contributions make the implementation
question concrete. They leave several questions on which the product's actual
banking value depends.

First, the product needs evidence about omissions across a \emph{complete}
regulatory packet, including incorporated definitions and attachments. A
benchmark whose input already supplies the decisive facts and formal rules
cannot answer whether a workbench discovers those facts and rules. The next
comparison should preserve independently inventoried source units and measure
which material units and alternatives are missed. Existing engineering cases
can remain as development checks without serving as the legal reference.

Second, comparison of language backends needs matched meanings. Before choosing
among RuleIR, Catala, a decision table or a normative language, give each the same
unknown facts, conflicting evidence, exceptions, dates and duties. Record where
an adapter cannot express the accepted meaning. A familiar syntax or a successful
example does not establish interchangeability. The restricted initial runtime
keeps this comparison manageable.

\begin{figure}[H]
\centering
\TeachingFlow{One reviewed banking question}{Same facts, authority and resource budget}{Compare consequential failures}
\caption{Coverage improves through matched questions, not by accumulating unrelated benchmark scores.}
\label{fig:literature-matched-question}
\end{figure}

Third, common-mode failure deserves a direct study. Changing the model vendor
while keeping the same defective extraction may leave the most serious omission
untouched. A useful experiment varies extraction, inventory and interpretation
separately and retains which activity first detected a defect. The numerical
dependence example in Chapter~\ref{ch:ensemble} explains the mechanism; it supplies
no measured correlation for this product.

Finally, the proposed human workflow needs independent evaluation. Reviewers
must be able to identify a plausible but unsupported source claim, notice a
minority reading and act on unresolved issues without an unmanageable queue.
The study in Chapter~\ref{ch:evaluation} therefore measures consequential errors
and review effort together. A legal-language theorem, model benchmark or
successful browser test does not answer that human question. These are defined
coverage gaps and research tasks, rather than a claim that no relevant work
exists beyond the retained literature.
